\section{Conclusion}
A scientific-verification result requires more than a score.  The intended
terminal must be identifiable from admissible scientific information, resistant
to nuisance decoding, expressible by the compared interfaces, and capable of
varying over the panel.  The exact-risk and factorization results make these
requirements testable: a target relation is implementable precisely when it is
constant on the decision representation's fibres and its terminals are
available; mixed fibres impose an irreducible lower bound.

The three P4 evidence layers instantiate different parts of that theory without
rewriting one another.  V2 eliminated false promotion across 360 hostile
opportunities while preserving 60/60 clean promotions, but its H3 family was
construction-saturated and remains \textsc{NotSupported}.  V3 repaired the
registered nuisance shortcuts and separated the exact \textsc{CannotCheck}
terminal, but its comparison is principally one of interface expressiveness.
P4-X then moved the comparator upward: after donor-complete provenance,
verification, custody, version, epoch, and generic-authorization mechanisms were
granted, the non-compensatory scientific-promotion relation scored 400/400
against 250/400 for the registered generic product and 50/400 for a compensatory
product, while an information-equivalent typed product tied at 400/400.  The tie
is a positive boundary result: scientific promotion need not be centralized;
it must be represented with target-sufficient state and the appropriate
relation.

The resulting contribution is wider than a gate checklist but narrower than a
deployed-system victory.  It is a theory of when verification axes can support
scientific claims, plus exact protected evidence that local verification and
generic permission need not suffice for scientific promotion.  The frozen
768-cluster naturalistic successor is the next discriminator and currently has
no external-panel outcome.  Naturalistic, verifier-provider, source-disjoint,
and universal scientific-authorization superiority therefore remain open
research targets, not promoted claims.

\bibliographystyle{tmlr}
\bibliography{bibliography}

\appendix
